\chapter{CIM Wizard Integrated Calculators}
\label{app:calculators}

The \texttt{cim\_wizard\_integrated} service exposes its modelling logic through a
collection of Python ``calculators'' hosted in
\texttt{app/calculators/}. Each class is initialised with the shared
\texttt{pipeline\_executor}, giving access to the FastAPI logger, a
stateful \texttt{data\_manager}, and---when available—an SQLAlchemy or
Django session. Calculators follow the same pattern:
\begin{enumerate}[label=\alph*)]
    \item read prerequisite features via \texttt{pipeline.get\_feature\_safely()};
    \item run a deterministic or data-driven method (often exposed as
    \texttt{calculate\_*} or \texttt{by\_*});
    \item store the result back into the pipeline context, typically under a
    feature name that matches the class (for example
    \texttt{building\_area});
    \item optionally persist the derived values into the PostGIS or 3DCityDB schema.
\end{enumerate}
The following sections describe each available calculator, grouped by
their role inside the CIM Wizard pipeline.

\section{Scenario and Census Initialisation}
\subsection{Scenario Geometry Calculator}
\textbf{Class:} \texttt{ScenarioGeoCalculator} (\texttt{scenario\_geo\_calculator.py})\\
\textbf{Main methods:}
\begin{itemize}[leftmargin=2em]
    \item \texttt{calculate\_from\_scenario\_geo()} – accepts the GeoJSON
    polygon submitted from the UI, generates missing project/scenario
    identifiers, computes the centroid if the user did not provide map
    centre information, and stores the derived boundary under
    \texttt{scenario\_geo}.
    \item \texttt{calculate\_from\_buildings\_geo()} – when only
    building footprints are available, the method unions them via
    Shapely/GeoPandas to obtain a convex hull that plays the role of the
    scenario boundary.
\end{itemize}
Both methods persist the boundary inside the Django model
\texttt{Project\_Scenario} when the ORM is available.

\subsection{Scenario Census Boundary Calculator}
\textbf{Class:} \texttt{ScenarioCensusBoundaryCalculator}\\
\textbf{Main method:} \texttt{calculate\_from\_census\_api()}\\
The class currently synthesises a census boundary by reusing the
scenario geometry and wrapping it inside a GeoJSON Feature that contains
one or more \texttt{census\_zones}. Each zone transports the attributes
needed by downstream calculators (E8–E16 construction counts, P1
population, PF1 families). When a SQLAlchemy session exists the method
stores the MultiPolygon in \texttt{project\_scenario.census\_boundary}.

\subsection{Census Population Calculator}
\textbf{Class:} \texttt{CensusPopulationCalculator}\\
\textbf{Main method:} \texttt{calculate\_from\_census\_boundary()}\\
This utility sums the \texttt{total\_population} field reported in
the \texttt{scenario\_census\_boundary} feature. When the aggregated
value is missing it falls back to the per-zone counts. The return value
is a scalar that becomes the global population budget for volume-based
distributions.

\section{Geometry Acquisition and Pre-processing}
\subsection{Building Geo Calculator}
\textbf{Class:} \texttt{BuildingGeoCalculator}\\
\textbf{Primary entry points:}
\begin{itemize}[leftmargin=2em]
    \item \texttt{calculate\_from\_scenario\_census\_geo()} – validates the
    scenario boundary, queries OpenStreetMap (via osmnx when available or
    via Overpass QL otherwise), filters the resulting polygons so that their
    centroids fall inside the project boundary, and returns a
    \texttt{FeatureCollection}-like payload under \texttt{building\_geo}.
    \item \texttt{calculate\_from\_building\_geo()} – ingests raw GeoJSON
    blobs submitted through the API, normalises IDs and properties, and
    stores them using the same schema as the OSM-based workflow.
\end{itemize}
Helper methods perform OSM usage classification, fallback height
estimation from tags, and database persistence of geometry records.

\subsection{Building Properties Initialiser}
\textbf{Class:} \texttt{BuildingPropsCalculator}\\
\textbf{Main method:} \texttt{init()}\\
Creates placeholder rows (height/area/volume/floors set to
\texttt{None}) for each footprint from \texttt{building\_geo}, writes
them to the vector schema (table \texttt{BuildingProperties}) when the
SQLAlchemy session is available, and exposes the list as
\texttt{building\_props}. Subsequent calculators progressively fill the
missing attributes.

\subsection{Building Geo LoD 1.2 Calculator}
\textbf{Class:} \texttt{BuildingGeoLod12Calculator}\\
\textbf{Main method:} \texttt{by\_footprint\_height()}\\
Generates simplified LoD~1.2 semantic surfaces when both footprints and
height estimates exist. For every building the calculator crafts
\texttt{WallSurface}, \texttt{RoofSurface}, and \texttt{GroundSurface}
records enriched with metadata (surface orientation, area, LoD tag) and
stores them inside the \texttt{building\_geo\_lod12} feature for later
export to 3DCityDB.

\section{Physical Feature Calculators}
\subsection{Area}
\textbf{Class:} \texttt{BuildingAreaCalculator}\\
\textbf{Method:} \texttt{calculate\_from\_geometry()}\\
Reads every footprint from \texttt{building\_geo}, computes the area
(Shapely projected to UTM~32N with an approximate fallback), logs the
five most recent samples, writes the results under both
\texttt{building\_area} and \texttt{building\_areas}, and pushes the
values into \texttt{BuildingProperties.area}.

\subsection{Height}
\textbf{Class:} \texttt{BuildingHeightCalculator}\\
\textbf{Method:} \texttt{calculate\_from\_raster\_tiles()}\\
Demonstrates the DSM--DTM approach: for the centroid of each building
the method queries \texttt{cim\_raster.dsm\_raster\_tiles} and
\texttt{cim\_wizard.dtm\_raster\_tiles} via SQLAlchemy, clamps the
result between 3~m and 200~m, and stores the list as
\texttt{building\_height}.

\subsection{Number of Floors}
\textbf{Class:} \texttt{BuildingNFloorsCalculator}\\
\textbf{Method:} \texttt{estimate\_by\_height()}\\
Consumes \texttt{building\_height} and the residential filter (see
Section~\ref{sec:resfilter}), divides each height by 3~m to obtain an
integer floor count, skips non-residential buildings, and updates both
the pipeline feature \texttt{building\_n\_floors} and
\texttt{BuildingProperties.number\_of\_floors}.

\subsection{Volume}
\textbf{Class:} \texttt{BuildingVolumeCalculator}\\
\textbf{Method:} \texttt{calculate\_from\_height\_and\_area()}\\
Crosses the previously computed heights, areas, and residential flags to
generate per-building volumes (height $\times$ area). The class logs
sample calculations, persists the values, and exposes them as both the
rich \texttt{building\_volume} feature and the convenience list
\texttt{building\_volumes}.

\section{Residential Filtering and Typology Assignment}
\subsection{Residential Filter}
\label{sec:resfilter}
\textbf{Class:} \texttt{BuildingResidentialFilterCalculator}\\
\textbf{Method:} \texttt{calculate\_filter\_res()}\\
Implements a conservative classifier that marks a footprint as
non-residential when \texttt{area}~$<90~\mathrm{m}^2$, height
${}<4~\mathrm{m}$, or the OSM \texttt{building}/\texttt{amenity} tags
match a curated non-residential list. The resulting boolean vector is
stored as \texttt{filter\_res} and statistics about filtering causes are
reported in the logs.

\subsection{Building Type Calculator}
\textbf{Class:} \texttt{BuildingTypeCalculator}\\
\textbf{Method:} \texttt{by\_census\_osm()}\\
Starting from the GeoDataFrames produced by the census integration
branch (Section~\ref{sec:demographics}), this class:
\begin{enumerate}[label=\roman*),leftmargin=2em]
    \item forces any building tagged as \texttt{not\_residential\_based\_on\_osm} to be classified as non-residential;
    \item applies strict thresholds (height $>$~8~m and area $>$~100~m$^2$) to
    the remaining footprints;
    \item records zone-level discrepancies between census counts and the
    assigned types (useful for accuracy reporting).
\end{enumerate}

\subsection{Construction Year Distribution}
\textbf{Class:} \texttt{BuildingConstructionYearCalculator}\\
\textbf{Method:} \texttt{by\_census\_osm()}\\
Maps the census E8--E16 buckets to year ranges and TABULA periods. For
each census zone, the function computes the percentage split,
generates random construction years inside every bucket, shuffles the
assignments to avoid bias, and attaches three attributes to every
residential building: \texttt{const\_period\_census},
\texttt{const\_year}, and \texttt{const\_TABULA}.

\section{Demographic Calculators}
\label{sec:demographics}
\subsection{Population Distribution}
\textbf{Class:} \texttt{BuildingPopulationCalculator}\\
\textbf{Methods:}
\begin{itemize}[leftmargin=2em]
    \item \texttt{calculate\_from\_volume\_distribution()} – uses the total
    population returned by the census calculator and distributes it
    proportionally to each residential building volume.
    \item \texttt{by\_census\_osm()} – operates on the GeoDataFrames,
    assigning \texttt{n\_people} by weighting the zone population with
    the ratio between the building volume and the total residential
    volume in that zone.
\end{itemize}

\subsection{Families per Building}
\textbf{Class:} \texttt{BuildingNFamiliesCalculator}\\
\textbf{Methods:}
\begin{itemize}[leftmargin=2em]
    \item \texttt{calculate\_from\_population()} – divides the per-building
    population by an average family size (2.5) to estimate
    \texttt{n\_family}.
    \item \texttt{by\_census\_osm()} – alternative version used inside the
    GeoDataFrame workflow (defaults to 3 persons per family).
\end{itemize}

\subsection{Building Demographic Calculator}
\textbf{Class:} \texttt{BuildingDemographicCalculator}\\
This orchestrator glues together the previous calculators when the
GeoDataFrame path is enabled:
\begin{enumerate}[label=\alph*),leftmargin=2em]
    \item fetches or calls the scenario census boundary service;
    \item assigns building types, construction years, volumes, population,
    and families by delegating to the specialised calculators;
    \item filters footprints to the project boundary and writes the enriched
    attributes back to the database;
    \item composes a final summary under \texttt{building\_demographic},
    including accuracy reports (for example, whether the number of
    residential buildings per zone matches census expectations).
\end{enumerate}

\section{Summary of Feature Dependencies}
Table~\ref{tab:calc-dependencies} summarises the inputs and outputs of
each calculator. The table can be used as a quick reference when adding
new methods to the pipeline.

\begin{table}[htbp]
    \centering
    \caption{Calculator inputs and produced features.}
    \label{tab:calc-dependencies}
    \begin{tabular}{p{3.6cm}p{5.2cm}p{5.8cm}}
        \toprule
        \textbf{Calculator} & \textbf{Required features / services} & \textbf{Produced features / side effects} \\
        \midrule
        ScenarioGeo & UI GeoJSON or building footprints & \texttt{scenario\_geo}, optional \texttt{Project\_Scenario} rows \\
        ScenarioCensusBoundary & \texttt{scenario\_geo} + census API & \texttt{scenario\_census\_boundary}, census boundary persisted to DB \\
        CensusPopulation & \texttt{scenario\_census\_boundary} & scalar population budget \\
        BuildingGeo & Scenario boundary or UI features & \texttt{building\_geo}, \texttt{Building} table records \\
        BuildingProps & \texttt{building\_geo} & \texttt{building\_props}, empty \texttt{BuildingProperties} rows \\
        BuildingArea & \texttt{building\_geo} & \texttt{building\_area}, \texttt{building\_areas}, DB areas \\
        BuildingHeight & \texttt{building\_geo}, raster tiles & \texttt{building\_height}, DB heights \\
        BuildingNFloors & \texttt{building\_height}, \texttt{filter\_res} & \texttt{building\_n\_floors}, DB floor counts \\
        BuildingVolume & \texttt{building\_height}, \texttt{building\_area}, \texttt{filter\_res} & \texttt{building\_volume}, \texttt{building\_volumes}, DB volumes \\
        BuildingResidentialFilter & \texttt{building\_geo}, \texttt{building\_area}, \texttt{building\_height} & \texttt{filter\_res} vector \\
        BuildingType & Census and building GeoDataFrames & Updated GeoDataFrame with \texttt{building\_type} + zone errors \\
        BuildingConstructionYear & Census and building GeoDataFrames & \texttt{const\_period\_census}, \texttt{const\_year}, \texttt{const\_TABULA} \\
        BuildingPopulation & \texttt{building\_volume} + census totals or GeoDataFrames & \texttt{building\_population}, \texttt{n\_people} per building \\
        BuildingNFamilies & \texttt{building\_population} & \texttt{building\_n\_families}, \texttt{n\_family} \\
        BuildingGeoLoD12 & \texttt{building\_geo}, \texttt{building\_height} & \texttt{building\_geo\_lod12} surfaces \\
        BuildingDemographic & All of the above GeoDataFrame-based features & \texttt{building\_demographic}, DB updates with demographic fields \\
        \bottomrule
    \end{tabular}
\end{table}

Collectively, these calculators demonstrate how CIM Wizard exposes an
extensible, object-oriented API for urban feature computation: each
module documents its expectations through the feature names and can be
replaced or expanded without touching the rest of the pipeline. This
appendix serves as the authoritative reference when integrating new
calculators or when tracing data provenance within the platform.

